\documentclass[11pt,a4paper]{article}

\newcommand{\doctitle}{Storm Sequence Generator (MKM-SS-001)}
\newcommand{\docsubtitle}{Wolfram Oracle Benchmark of the Code and its v3.1 Document}
\newcommand{\docversion}{0.3}
\newcommand{\docdate}{09-November-2026}
\newcommand{\docauthor}{David K Kelly}

% MKM Research Labs — Shared Document Header
% Usage: % MKM Research Labs — Shared Document Header
% Usage: % MKM Research Labs — Shared Document Header
% Usage: \input{../shared/mkm_header}  (before \begin{document})
%
% Required commands before \input:
%   \newcommand{\doctitle}{...}
%   \newcommand{\docsubtitle}{...}
%   \newcommand{\docversion}{1.0}
%   \newcommand{\docdate}{13-February-2026}
%   \newcommand{\docauthor}{David K Kelly}

% ─── Packages ───────────────────────────────────────────────────────────────
\usepackage[utf8]{inputenc}
\usepackage{amsmath,amsfonts,amssymb}
\usepackage{listings}
\usepackage{color,xcolor}
\usepackage{graphicx}
\usepackage{hyperref}
\usepackage[toc,page]{appendix}
\usepackage{tabto}
\usepackage{float}
\usepackage{comment}
\usepackage[en-GB]{datetime2}
\usepackage{setspace}
\usepackage{subcaption}
\usepackage{booktabs}
\usepackage{longtable}
\usepackage{adjustbox}
\usepackage{enumitem}
\usepackage{fancyhdr}
\usepackage{lastpage}

% ─── Page geometry ──────────────────────────────────────────────────────────
\onehalfspacing
\usepackage{geometry}
\geometry{a4paper, total={170mm,247mm}, left=20mm, top=25mm, bottom=25mm}

% ─── Headers and footers ───────────────────────────────────────────────────
\pagestyle{fancy}
\fancyhf{}
\fancyhead[L]{\small\textcolor{gray}{MKM Research Labs}}
\fancyhead[R]{\small\textcolor{gray}{\doctitle{} --- v\docversion}}
\fancyfoot[C]{\small\textcolor{gray}{CONFIDENTIAL}}
\fancyfoot[R]{\small\textcolor{gray}{Page \thepage\ of \pageref{LastPage}}}
\renewcommand{\headrulewidth}{0.4pt}
\renewcommand{\footrulewidth}{0.4pt}

% ─── Colours ────────────────────────────────────────────────────────────────
\definecolor{mkmblue}{HTML}{1565C0}
\definecolor{codegreen}{rgb}{0,0.6,0}
\definecolor{codegray}{rgb}{0.5,0.5,0.5}
\definecolor{codepurple}{rgb}{0.58,0,0.82}
\definecolor{backcolour}{rgb}{0.95,0.95,0.92}

% ─── Hyperref setup ─────────────────────────────────────────────────────────
\hypersetup{
    colorlinks=true,
    linkcolor=mkmblue,
    citecolor=mkmblue,
    urlcolor=mkmblue
}

% ─── Code listings ──────────────────────────────────────────────────────────
\lstdefinestyle{mystyle}{
    backgroundcolor=\color{backcolour},
    commentstyle=\color{codegreen},
    keywordstyle=\color{magenta},
    numberstyle=\tiny\color{codegray},
    stringstyle=\color{codepurple},
    basicstyle=\ttfamily\footnotesize,
    breakatwhitespace=false,
    breaklines=true,
    captionpos=b,
    keepspaces=true,
    numbers=left,
    numbersep=5pt,
    showspaces=false,
    showstringspaces=false,
    showtabs=false,
    tabsize=2
}
\lstset{style=mystyle}

% ─── Section formatting ────────────────────────────────────────────────────
\setcounter{tocdepth}{3}
\setcounter{secnumdepth}{3}

% ─── Shared macros ──────────────────────────────────────────────────────────
\newcommand{\mkmtitlepage}{%
    \begin{titlepage}
    \begin{center}
    \vspace*{1cm}
    \begin{figure}[H]\centering
        \includegraphics[width=0.4\textwidth]{../shared/Logos/MKM_logo.png}
    \end{figure}
    \vspace{1.5cm}
    {\Huge\bfseries\textcolor{mkmblue}{\doctitle}\par}
    \vspace{0.5cm}
    {\Large\textcolor{gray}{\docsubtitle}\par}
    \vspace{2cm}
    {\large MKM Research Labs\par}
    \vspace{0.5cm}
    {\large Version \docversion{} --- \docdate\par}
    \vspace{0.5cm}
    {\large \docauthor\par}
    \vspace{2cm}
    \textcolor{red}{\textbf{CONFIDENTIAL --- PROPRIETARY}}
    \end{center}
    \end{titlepage}
}

\newcommand{\mkmlegalpage}{%
    \clearpage
    \thispagestyle{empty}
    \begin{center}
    \large\textbf{Legal Notice}
    \vspace{1em}

    \textbf{PROPRIETARY AND CONFIDENTIAL}
    \end{center}

    \noindent This document, including all algorithms, models, methodology, formulas,
    calculations, and intellectual concepts contained herein, constitutes the exclusive
    intellectual property and confidential trade secrets of MKM Research Labs (``MKM'').

    \vspace{0.5em}
    \noindent Any usage, reproduction, distribution, modification, or disclosure of this
    document or any portion thereof, without the express prior written authorisation from
    MKM Research Labs is strictly prohibited and constitutes an infringement of intellectual
    property rights.

    \vspace{0.5em}
    \noindent The document is provided on an ``as is'' basis. MKM Research Labs makes no
    representations or warranties of any kind, express or implied, regarding its accuracy,
    reliability, or suitability for any particular purpose.

    \vspace{1em}
    \begin{center}
    All rights reserved. \copyright\ 2021--2026 MKM Research Labs.

    \vspace{0.5em}
    Governed by the laws of the United Kingdom.
    \end{center}
    \clearpage
}

\newcommand{\mkmhistorypage}[1]{%
    \section{Document History}
    \begin{table}[ht]
        \centering
        \begin{tabular}{l|l|l|l}
            \textbf{Date} & \textbf{Version} & \textbf{Description} & \textbf{Author} \\
            \hline
            #1
        \end{tabular}
        \label{tab:revision_history}
    \end{table}
    \clearpage
}

% ─── Bank validation protocol section commands ──────────────────────────────
% Use these to ensure every model doc has the required sections.
% They are just styled \section calls — content is provided by each doc.

\newcommand{\mkmexecsummary}{\section{Executive Summary}}
\newcommand{\mkmpurpose}{\section{Model Purpose and Scope}}
\newcommand{\mkmframework}{\section{Mathematical Framework}}
\newcommand{\mkminputs}{\section{Input Parameters and Calibration}}
\newcommand{\mkmimplementation}{\section{Implementation Details}}
\newcommand{\mkmsensitivity}{\section{Sensitivity Analysis}}
\newcommand{\mkmvalidation}{\section{Validation and Backtesting}}
\newcommand{\mkmtestresults}{\subsection{Automated Test Results}\input{test_results}}
\newcommand{\mkmlimitations}{\section{Model Limitations and Known Weaknesses}}

\newcommand{\mkmgovernance}{\section{Governance Validation Questions}
\noindent The following nine questions are mandated by Chapter~5 of the
\emph{Model Risk Governance for Vendors} handbook and must be explicitly
addressed for every model in the inventory. Response status is tracked in the
Model Governance Dashboard (MRC portal). \par\medskip}

\newcommand{\mkmvqitem}[3]{\subsection{Q#1 --- #2}
\noindent\textit{``#3''}\par\medskip\noindent}

\newcommand{\mkmchangelog}{\section{Change History}}
  (before \begin{document})
%
% Required commands before \input:
%   \newcommand{\doctitle}{...}
%   \newcommand{\docsubtitle}{...}
%   \newcommand{\docversion}{1.0}
%   \newcommand{\docdate}{13-February-2026}
%   \newcommand{\docauthor}{David K Kelly}

% ─── Packages ───────────────────────────────────────────────────────────────
\usepackage[utf8]{inputenc}
\usepackage{amsmath,amsfonts,amssymb}
\usepackage{listings}
\usepackage{color,xcolor}
\usepackage{graphicx}
\usepackage{hyperref}
\usepackage[toc,page]{appendix}
\usepackage{tabto}
\usepackage{float}
\usepackage{comment}
\usepackage[en-GB]{datetime2}
\usepackage{setspace}
\usepackage{subcaption}
\usepackage{booktabs}
\usepackage{longtable}
\usepackage{adjustbox}
\usepackage{enumitem}
\usepackage{fancyhdr}
\usepackage{lastpage}

% ─── Page geometry ──────────────────────────────────────────────────────────
\onehalfspacing
\usepackage{geometry}
\geometry{a4paper, total={170mm,247mm}, left=20mm, top=25mm, bottom=25mm}

% ─── Headers and footers ───────────────────────────────────────────────────
\pagestyle{fancy}
\fancyhf{}
\fancyhead[L]{\small\textcolor{gray}{MKM Research Labs}}
\fancyhead[R]{\small\textcolor{gray}{\doctitle{} --- v\docversion}}
\fancyfoot[C]{\small\textcolor{gray}{CONFIDENTIAL}}
\fancyfoot[R]{\small\textcolor{gray}{Page \thepage\ of \pageref{LastPage}}}
\renewcommand{\headrulewidth}{0.4pt}
\renewcommand{\footrulewidth}{0.4pt}

% ─── Colours ────────────────────────────────────────────────────────────────
\definecolor{mkmblue}{HTML}{1565C0}
\definecolor{codegreen}{rgb}{0,0.6,0}
\definecolor{codegray}{rgb}{0.5,0.5,0.5}
\definecolor{codepurple}{rgb}{0.58,0,0.82}
\definecolor{backcolour}{rgb}{0.95,0.95,0.92}

% ─── Hyperref setup ─────────────────────────────────────────────────────────
\hypersetup{
    colorlinks=true,
    linkcolor=mkmblue,
    citecolor=mkmblue,
    urlcolor=mkmblue
}

% ─── Code listings ──────────────────────────────────────────────────────────
\lstdefinestyle{mystyle}{
    backgroundcolor=\color{backcolour},
    commentstyle=\color{codegreen},
    keywordstyle=\color{magenta},
    numberstyle=\tiny\color{codegray},
    stringstyle=\color{codepurple},
    basicstyle=\ttfamily\footnotesize,
    breakatwhitespace=false,
    breaklines=true,
    captionpos=b,
    keepspaces=true,
    numbers=left,
    numbersep=5pt,
    showspaces=false,
    showstringspaces=false,
    showtabs=false,
    tabsize=2
}
\lstset{style=mystyle}

% ─── Section formatting ────────────────────────────────────────────────────
\setcounter{tocdepth}{3}
\setcounter{secnumdepth}{3}

% ─── Shared macros ──────────────────────────────────────────────────────────
\newcommand{\mkmtitlepage}{%
    \begin{titlepage}
    \begin{center}
    \vspace*{1cm}
    \begin{figure}[H]\centering
        \includegraphics[width=0.4\textwidth]{../shared/Logos/MKM_logo.png}
    \end{figure}
    \vspace{1.5cm}
    {\Huge\bfseries\textcolor{mkmblue}{\doctitle}\par}
    \vspace{0.5cm}
    {\Large\textcolor{gray}{\docsubtitle}\par}
    \vspace{2cm}
    {\large MKM Research Labs\par}
    \vspace{0.5cm}
    {\large Version \docversion{} --- \docdate\par}
    \vspace{0.5cm}
    {\large \docauthor\par}
    \vspace{2cm}
    \textcolor{red}{\textbf{CONFIDENTIAL --- PROPRIETARY}}
    \end{center}
    \end{titlepage}
}

\newcommand{\mkmlegalpage}{%
    \clearpage
    \thispagestyle{empty}
    \begin{center}
    \large\textbf{Legal Notice}
    \vspace{1em}

    \textbf{PROPRIETARY AND CONFIDENTIAL}
    \end{center}

    \noindent This document, including all algorithms, models, methodology, formulas,
    calculations, and intellectual concepts contained herein, constitutes the exclusive
    intellectual property and confidential trade secrets of MKM Research Labs (``MKM'').

    \vspace{0.5em}
    \noindent Any usage, reproduction, distribution, modification, or disclosure of this
    document or any portion thereof, without the express prior written authorisation from
    MKM Research Labs is strictly prohibited and constitutes an infringement of intellectual
    property rights.

    \vspace{0.5em}
    \noindent The document is provided on an ``as is'' basis. MKM Research Labs makes no
    representations or warranties of any kind, express or implied, regarding its accuracy,
    reliability, or suitability for any particular purpose.

    \vspace{1em}
    \begin{center}
    All rights reserved. \copyright\ 2021--2026 MKM Research Labs.

    \vspace{0.5em}
    Governed by the laws of the United Kingdom.
    \end{center}
    \clearpage
}

\newcommand{\mkmhistorypage}[1]{%
    \section{Document History}
    \begin{table}[ht]
        \centering
        \begin{tabular}{l|l|l|l}
            \textbf{Date} & \textbf{Version} & \textbf{Description} & \textbf{Author} \\
            \hline
            #1
        \end{tabular}
        \label{tab:revision_history}
    \end{table}
    \clearpage
}

% ─── Bank validation protocol section commands ──────────────────────────────
% Use these to ensure every model doc has the required sections.
% They are just styled \section calls — content is provided by each doc.

\newcommand{\mkmexecsummary}{\section{Executive Summary}}
\newcommand{\mkmpurpose}{\section{Model Purpose and Scope}}
\newcommand{\mkmframework}{\section{Mathematical Framework}}
\newcommand{\mkminputs}{\section{Input Parameters and Calibration}}
\newcommand{\mkmimplementation}{\section{Implementation Details}}
\newcommand{\mkmsensitivity}{\section{Sensitivity Analysis}}
\newcommand{\mkmvalidation}{\section{Validation and Backtesting}}
\newcommand{\mkmtestresults}{\subsection{Automated Test Results}\input{test_results}}
\newcommand{\mkmlimitations}{\section{Model Limitations and Known Weaknesses}}

\newcommand{\mkmgovernance}{\section{Governance Validation Questions}
\noindent The following nine questions are mandated by Chapter~5 of the
\emph{Model Risk Governance for Vendors} handbook and must be explicitly
addressed for every model in the inventory. Response status is tracked in the
Model Governance Dashboard (MRC portal). \par\medskip}

\newcommand{\mkmvqitem}[3]{\subsection{Q#1 --- #2}
\noindent\textit{``#3''}\par\medskip\noindent}

\newcommand{\mkmchangelog}{\section{Change History}}
  (before \begin{document})
%
% Required commands before \input:
%   \newcommand{\doctitle}{...}
%   \newcommand{\docsubtitle}{...}
%   \newcommand{\docversion}{1.0}
%   \newcommand{\docdate}{13-February-2026}
%   \newcommand{\docauthor}{David K Kelly}

% ─── Packages ───────────────────────────────────────────────────────────────
\usepackage[utf8]{inputenc}
\usepackage{amsmath,amsfonts,amssymb}
\usepackage{listings}
\usepackage{color,xcolor}
\usepackage{graphicx}
\usepackage{hyperref}
\usepackage[toc,page]{appendix}
\usepackage{tabto}
\usepackage{float}
\usepackage{comment}
\usepackage[en-GB]{datetime2}
\usepackage{setspace}
\usepackage{subcaption}
\usepackage{booktabs}
\usepackage{longtable}
\usepackage{adjustbox}
\usepackage{enumitem}
\usepackage{fancyhdr}
\usepackage{lastpage}

% ─── Page geometry ──────────────────────────────────────────────────────────
\onehalfspacing
\usepackage{geometry}
\geometry{a4paper, total={170mm,247mm}, left=20mm, top=25mm, bottom=25mm}

% ─── Headers and footers ───────────────────────────────────────────────────
\pagestyle{fancy}
\fancyhf{}
\fancyhead[L]{\small\textcolor{gray}{MKM Research Labs}}
\fancyhead[R]{\small\textcolor{gray}{\doctitle{} --- v\docversion}}
\fancyfoot[C]{\small\textcolor{gray}{CONFIDENTIAL}}
\fancyfoot[R]{\small\textcolor{gray}{Page \thepage\ of \pageref{LastPage}}}
\renewcommand{\headrulewidth}{0.4pt}
\renewcommand{\footrulewidth}{0.4pt}

% ─── Colours ────────────────────────────────────────────────────────────────
\definecolor{mkmblue}{HTML}{1565C0}
\definecolor{codegreen}{rgb}{0,0.6,0}
\definecolor{codegray}{rgb}{0.5,0.5,0.5}
\definecolor{codepurple}{rgb}{0.58,0,0.82}
\definecolor{backcolour}{rgb}{0.95,0.95,0.92}

% ─── Hyperref setup ─────────────────────────────────────────────────────────
\hypersetup{
    colorlinks=true,
    linkcolor=mkmblue,
    citecolor=mkmblue,
    urlcolor=mkmblue
}

% ─── Code listings ──────────────────────────────────────────────────────────
\lstdefinestyle{mystyle}{
    backgroundcolor=\color{backcolour},
    commentstyle=\color{codegreen},
    keywordstyle=\color{magenta},
    numberstyle=\tiny\color{codegray},
    stringstyle=\color{codepurple},
    basicstyle=\ttfamily\footnotesize,
    breakatwhitespace=false,
    breaklines=true,
    captionpos=b,
    keepspaces=true,
    numbers=left,
    numbersep=5pt,
    showspaces=false,
    showstringspaces=false,
    showtabs=false,
    tabsize=2
}
\lstset{style=mystyle}

% ─── Section formatting ────────────────────────────────────────────────────
\setcounter{tocdepth}{3}
\setcounter{secnumdepth}{3}

% ─── Shared macros ──────────────────────────────────────────────────────────
\newcommand{\mkmtitlepage}{%
    \begin{titlepage}
    \begin{center}
    \vspace*{1cm}
    \begin{figure}[H]\centering
        \includegraphics[width=0.4\textwidth]{../shared/Logos/MKM_logo.png}
    \end{figure}
    \vspace{1.5cm}
    {\Huge\bfseries\textcolor{mkmblue}{\doctitle}\par}
    \vspace{0.5cm}
    {\Large\textcolor{gray}{\docsubtitle}\par}
    \vspace{2cm}
    {\large MKM Research Labs\par}
    \vspace{0.5cm}
    {\large Version \docversion{} --- \docdate\par}
    \vspace{0.5cm}
    {\large \docauthor\par}
    \vspace{2cm}
    \textcolor{red}{\textbf{CONFIDENTIAL --- PROPRIETARY}}
    \end{center}
    \end{titlepage}
}

\newcommand{\mkmlegalpage}{%
    \clearpage
    \thispagestyle{empty}
    \begin{center}
    \large\textbf{Legal Notice}
    \vspace{1em}

    \textbf{PROPRIETARY AND CONFIDENTIAL}
    \end{center}

    \noindent This document, including all algorithms, models, methodology, formulas,
    calculations, and intellectual concepts contained herein, constitutes the exclusive
    intellectual property and confidential trade secrets of MKM Research Labs (``MKM'').

    \vspace{0.5em}
    \noindent Any usage, reproduction, distribution, modification, or disclosure of this
    document or any portion thereof, without the express prior written authorisation from
    MKM Research Labs is strictly prohibited and constitutes an infringement of intellectual
    property rights.

    \vspace{0.5em}
    \noindent The document is provided on an ``as is'' basis. MKM Research Labs makes no
    representations or warranties of any kind, express or implied, regarding its accuracy,
    reliability, or suitability for any particular purpose.

    \vspace{1em}
    \begin{center}
    All rights reserved. \copyright\ 2021--2026 MKM Research Labs.

    \vspace{0.5em}
    Governed by the laws of the United Kingdom.
    \end{center}
    \clearpage
}

\newcommand{\mkmhistorypage}[1]{%
    \section{Document History}
    \begin{table}[ht]
        \centering
        \begin{tabular}{l|l|l|l}
            \textbf{Date} & \textbf{Version} & \textbf{Description} & \textbf{Author} \\
            \hline
            #1
        \end{tabular}
        \label{tab:revision_history}
    \end{table}
    \clearpage
}

% ─── Bank validation protocol section commands ──────────────────────────────
% Use these to ensure every model doc has the required sections.
% They are just styled \section calls — content is provided by each doc.

\newcommand{\mkmexecsummary}{\section{Executive Summary}}
\newcommand{\mkmpurpose}{\section{Model Purpose and Scope}}
\newcommand{\mkmframework}{\section{Mathematical Framework}}
\newcommand{\mkminputs}{\section{Input Parameters and Calibration}}
\newcommand{\mkmimplementation}{\section{Implementation Details}}
\newcommand{\mkmsensitivity}{\section{Sensitivity Analysis}}
\newcommand{\mkmvalidation}{\section{Validation and Backtesting}}
\newcommand{\mkmtestresults}{\subsection{Automated Test Results}\input{test_results}}
\newcommand{\mkmlimitations}{\section{Model Limitations and Known Weaknesses}}

\newcommand{\mkmgovernance}{\section{Governance Validation Questions}
\noindent The following nine questions are mandated by Chapter~5 of the
\emph{Model Risk Governance for Vendors} handbook and must be explicitly
addressed for every model in the inventory. Response status is tracked in the
Model Governance Dashboard (MRC portal). \par\medskip}

\newcommand{\mkmvqitem}[3]{\subsection{Q#1 --- #2}
\noindent\textit{``#3''}\par\medskip\noindent}

\newcommand{\mkmchangelog}{\section{Change History}}

% T1 so the underscores in code identifiers print as underscores.
\usepackage[T1]{fontenc}
\usepackage{lmodern}
\usepackage{longtable}
% Generated by tools.wolfram_oracle -- do not edit.
\newcommand{\bmGenerated}{2026-11-09}
\newcommand{\bmCodeCommit}{\texttt{28ef2f7f5745}}
\newcommand{\bmCodeDirty}{ (with uncommitted changes)}
\newcommand{\bmNumerics}{numpy 1.26.4, scipy 1.15.1, mpmath 1.4.1}
\newcommand{\bmWolfram}{15.0.0 for Mac OS X ARM (64-bit) (May 26, 2026)}
\newcommand{\bmWolframVia}{ via live kernel session (wolframclient)}
\newcommand{\bmWolframConfirms}{Wolfram, computing each probability independently, agrees with the reference to $1.9\times10^{-31}$.}
\newcommand{\bmChecks}{8}
\newcommand{\bmOutcome}{all pass}
\newcommand{\bmReplayN}{19{,}600}
\newcommand{\bmCodeErr}{$2.9\times10^{-16}$}
\newcommand{\bmPinnedPct}{3.1\%}
\newcommand{\bmCapPct}{8.6\%}
\newcommand{\bmShortPct}{35.4\%}
\newcommand{\bmSubsequent}{0.85}
\newcommand{\bmSeasonMonths}{1}
\newcommand{\bmIsolated}{5{,}225}
\newcommand{\bmClaims}{18}
\newcommand{\bmClaimsMet}{18}
\newcommand{\bmClaimsFailed}{0}

% The shared history table sets every column to its natural width; wrap the description.
\renewcommand{\mkmhistorypage}[1]{%
    \section{Document History}
    \begin{table}[ht]
        \centering
        \begin{tabular}{l|l|p{9cm}|l}
            \textbf{Date} & \textbf{Version} & \textbf{Description} & \textbf{Author} \\
            \hline
            #1
        \end{tabular}
    \end{table}
}

\begin{document}
% Long code paths and identifiers cannot hyphenate; let the spacing stretch instead.
\sloppy

\mkmtitlepage
\mkmlegalpage

\tableofcontents
\clearpage

% =========================================================================
\mkmexecsummary
% =========================================================================

This document benchmarks the Storm Sequence Generator (MKM-SS-001): the generators in
\texttt{src/port/src/storm\_multi/}, which draw the 10{,}000 multi-storm sequences the
stress pipeline turns into hydrographs and trains its flood classifiers on. It asks two
questions. Does the code compute what its model document (\texttt{storm\_multi\_v3.1.pdf},
v3.1, 9 November 2026) says? And does what the document says hold?

\paragraph{The code computes its documented equations.} A reference written from the
document's text replays \bmReplayN{} generated sequences on the code's own random stream.
Every sequence type and storm count matches, and every duration, gap, start time,
intensity and aggregate agrees to \bmCodeErr. In bulk, the code's sequence types, storm
counts and capped sequences follow the exact law the document's tables imply, within
sampling error (Table~\ref{tab:bm-checks}). \bmWolframConfirms{}

\paragraph{The document's claims hold:} \bmClaimsMet{} of \bmClaims{}
(Table~\ref{tab:bm-claims}).
\begin{itemize}[leftmargin=*,nosep]
\item \textbf{The recalibrated tables are the code's.} Table~1's storm counts (a cluster
  has three, a persistent sequence four) and Table~3's durations, a third of v3.0's, are
  what \texttt{config.port} holds, and the code draws from them.
\item \textbf{The storms fit the window.} The cap binds on \bmCapPct{} of duration draws
  (v3.0: 60.5\%). Gaps are shortened in \bmShortPct{} of multi-storm sequences. Every storm
  is capped in \bmPinnedPct{}. Both of Q6's thresholds hold at the baseline, as v3.1 says.
\item \textbf{Every gap is at least its minimum,} and Table~9's counts (\bmIsolated{}
  isolated and so on) are the exact expectation.
\item \textbf{The correlation is as stated:} a later storm is at least as strong as the
  first with probability \bmSubsequent.
\item \textbf{The override reaches the generator.} \texttt{storm\_control.json} changes the
  sequence probability and the duration table whether it is applied before or after the
  generator is imported.
\item \textbf{The paths, command line, outputs and modules} are those in the repository,
  and \textbf{the classifier figures} are those of the classifiers retrained on v3.1's
  catalogue.
\end{itemize}

\paragraph{What the benchmark does not settle.} v3.1's tables are calibrated to the
window, not to observed storms, and the benchmark cannot say whether they are realistic.
Each sequence's \texttt{sequence\_start}, which the stress pipeline reads as the season, is
the time it was generated (the default batch here: \bmSeasonMonths{} distinct month). The
document records both as limitations; neither is a claim this benchmark could pass or
fail.

\paragraph{Evidence.} \bmChecks{} checks, outcome: \bmOutcome{}
(Table~\ref{tab:bm-checks}).
\begin{itemize}[leftmargin=*,nosep]
\item Wolfram: \bmWolfram{}\bmWolframVia.
\item Code under test, at commit \bmCodeCommit{}\bmCodeDirty.
\item Numerics: \bmNumerics.
\item Generated \bmGenerated{} by \texttt{python -m tools.wolfram\_oracle ss001}.
\end{itemize}

% =========================================================================
\mkmpurpose
% =========================================================================

\subsection{What is benchmarked}
The generation path the document describes (\S5.1--\S5.2), loaded as its own modules:
\begin{itemize}[leftmargin=*,nosep]
\item \texttt{generators/batch\_generator.py}: \texttt{generate\_event\_set}, the category and
  base-intensity draws;
\item \texttt{generators/sequence\_generator.py}: type, count, the cap with its reserved gap
  minima, gap shortening, placement, intensities, precipitation and the aggregates;
\item the duration, gap and intensity samplers, and the tables in \texttt{config.port};
\item \texttt{utils/validation.py}.
\end{itemize}
Also tested: the repository the document cites, namely its paths, command line, outputs
and override control, and the classifier evidence.

\subsection{What is not}
\begin{itemize}[leftmargin=*,nosep]
\item The package's \texttt{models/} subpackage (sequence response, spatial correlation,
  hydrology), which the document places outside the model.
\item Whether the sequence types, durations and gaps are realistic. The document cites no
  observed sequences, and none are tested.
\item The classifier itself (MKM-FC-001). Only the figures this document cites from its
  training summary are checked.
\end{itemize}
This is a benchmark of the implementation (validation question Q8), not independent
validation.

\subsection{Earlier issue}
Benchmark v0.1 (7 November 2026) tested v2.0 and found that 11 of its 17 claims did not
hold. It is kept as evidence of the check. v3.0 is the model and document that answered
it. Benchmark v0.2 (8 November 2026) tested v3.0 and found all 17 of its claims held,
including its own record that the cap bound on 60.5\% of draws, above its Q6 threshold.
v3.1 recalibrates Tables~1 and~3 to resolve that, and this issue tests v3.1.

% =========================================================================
\mkmframework
% =========================================================================

The reference implements the document's equations from its text, in two halves.

\paragraph{Exact.} In rational arithmetic:
\begin{itemize}[leftmargin=*,nosep]
\item the type mix of Eq.~4 and Table~2, weighted by Table~6;
\item the storm counts of Table~1;
\item the chance that Eq.~6's cap, $(W_P - (n-1)\,g_{\min})/n$, binds;
\item the chance that every storm is capped, which holds every gap at its minimum;
\item the chance that Eq.~7 triggers Eq.~8's shortening;
\item the chance that a later storm is at least as strong as the first (Eqs.~12--13).
\end{itemize}
Durations and gaps are rounded to whole hours, as the document says.

\paragraph{Replay.} Eqs.~4--20 re-run on a \texttt{RandomState} seeded as the code's, drawing
in the code's order. Timings are exact rationals, so ``every storm capped'' is decided
without rounding; intensities are worked to 50 digits. The first storm's redraw (\S3.7) is
mirrored so the two streams stay in step.

% =========================================================================
\section{Benchmark Method}
% =========================================================================

\paragraph{The cases.}
\begin{itemize}[leftmargin=*,nosep]
\item The default batch: 10{,}000 sequences, seed 42, Table~6's weights, the Thames.
\item Every category with every forced type: 400 sequences each, 9{,}600 in all.
\item 100 seeds of 10{,}000 for the type mix (\S7.2).
\end{itemize}

\paragraph{The checks.}
\begin{itemize}[leftmargin=*,nosep]
\item The replay compares each sequence number for number.
\item The Monte Carlo compares the code's shares with the exact ones, in standard errors.
\item The validator runs on every sequence generated, and on a sequence broken five ways.
\end{itemize}

\paragraph{The statements.}
\begin{itemize}[leftmargin=*,nosep]
\item Figures (Table~9, \S6.2, \S6.3, \S7.2, Table~8) and Q6's baseline rates are judged
  against the exact law, to the precision the document states them.
\item Every gap is judged against its type's minimum, over all \bmReplayN{} sequences.
\item Q4's override is judged by applying a \texttt{storm\_control.json} that switches
  sequences off and shortens every storm. This runs in a separate interpreter, once with
  the generator imported before the file is applied and once after.
\item Paths, command-line flags and files are judged by searching the repository.
\item The classifier figures are judged against the training summary in the data store.
\end{itemize}

\paragraph{Wolfram.} Wolfram works from its own \texttt{TriangularDistribution} and
\texttt{UniformDistribution}, independently of the reference:
\begin{itemize}[leftmargin=*,nosep]
\item the type mix;
\item the cap's chance of binding in each configuration, as a single tail of the
  triangular law (the reference sums the rounded law);
\item the chance that every storm is capped, and of shortening, by convolving the rounded
  laws;
\item the correlation, by \texttt{Probability} over Eq.~13's mixture.
\end{itemize}

% =========================================================================
\mkmvalidation
% =========================================================================

% Generated by tools.wolfram_oracle -- do not edit.
\begin{table}[H]\centering\small
\caption{Benchmark checks}\label{tab:bm-checks}
\begin{adjustbox}{max width=\textwidth}
\begin{tabular}{p{96mm}p{16mm}rl}
\toprule
Check & Result & Limit & Outcome \\
\midrule
Replay, 19,600 sequences: types or storm counts that differ from the reference's & 0 & 0 & pass \\
Replay: worst relative error over durations, gaps, start times, intensities, precipitation, peak positions and Eqs.~16--20 & $2.9\times10^{-16}$ & $1.0\times10^{-12}$ & pass \\
Type mix over 100 seeds of 10{,}000 against Eq.~4 and Table~2's exact law (standard errors) & 0.37 & 4 & pass \\
Storm counts in the default batch against the exact law (standard errors) & 1.66 & 4 & pass \\
Multi-storm sequences with every storm capped, against the exact share (standard errors) & 0.00 & 4 & pass \\
\S5.4, the validator on all 19,600 sequences generated: invalid & 0 & 0 & pass \\
\S5.4, the validator: injected faults caught (late end, overlap, gap count, type, short gap) & 5 of 5 & 5 of 5 & pass \\
Wolfram agrees with the reference & $1.9\times10^{-31}$ & $1.0\times10^{-25}$ & pass \\
\bottomrule
\end{tabular}
\end{adjustbox}
\end{table}


\subsection{The model document's claims}
% Generated by tools.wolfram_oracle -- do not edit.
{\footnotesize\setlength{\tabcolsep}{4pt}\begin{longtable}{p{19mm}p{46mm}p{68mm}l}
\caption{The model document's claims, tested}\label{tab:bm-claims}\\
\toprule
Claim & Stated & Benchmark finds & Outcome \\
\midrule
\endhead
Eq.~4, Tables~2--6, 8 & the code's parameters & all as stated & met \\
Eq.~5, \S3.5 & durations and gaps are triangular draws & rounded to whole hours: 10,229 of 10,229 uncapped durations are whole; gaps are rounded before Eq.~8 scales them & \textbf{not met} \\
\S3.1, \S3.4.1, \S3.5.1 & every sequence ends by 156\,h, leaving 12\,h to drain & latest end 156.0\,h; least drainage 12.0\,h & met \\
\S3.5, Table~4 & gaps of 6--36\,h or 24--72\,h let the catchment drain between storms & 45.6\% of multi-storm sequences have no gap at all (100\% of catastrophic persistent ones); 82.1\% have gaps shorter than drawn & \textbf{not met} \\
Table~8 & a later storm is at least as strong as the first with probability 0.70 & 0.85 exactly; the code 0.853 over 10,881 storms. 0.70 is Eq.~13's chance of exceeding the base & \textbf{not met} \\
Table~7 & $P_{\mathrm{base}}$: Thames 35, Rhine 40, Danube 45, Mississippi 55\,mm & the code: 35, 35, 35, 35 & \textbf{not met} \\
\S6.1, Table~9 & expected counts 5,800, 2,000, 1,600, 600 (isolated, doublet, cluster, persistent) & exactly 5,225, 1,469, 2,116, 1,190 & \textbf{not met} \\
\S7.2 & over 100 seeds, every type within 2\% of Table~9 & worst 110\% off & \textbf{not met} \\
\S6.2 & the cap binds on $\approx$99\% of catastrophic durations at $n=5$, and $\approx$35\% of all draws & 100\%, and 50.7\% & \textbf{not met} \\
\S6.3 & gaps are scaled in $\approx$20\% of moderate-to-severe persistent sequences, $\approx$60\% of catastrophic ones & shorter than drawn in 99.2\% and 100\% & \textbf{not met} \\
\S5.3 & a fixed seed reproduces the batch & it does & met \\
\S5.4, \S7.1 & the validator finds no invalid sequence & 0 of 19,600; it catches 4 of 4 injected faults & met \\
Q2 & 10{,}000 sequences in under 60\,s & 0.4\,s & met \\
\S5.1 & the package holds the modules and functions listed & all present & met \\
\S2.2, \S5.5 & the consumers' paths, the command line and its output files & do not exist: \texttt{src/port/src/stressm.py}, \texttt{src/routes/trading/port\_stress.py}, \texttt{app.py}, \texttt{-{}-num-storms}, \texttt{-{}-train-classifier}, \texttt{doublet\_sequences.json} & \textbf{not met} \\
Q4 & no mechanism exists for runtime override of the tables & \texttt{storm\_control.json} overrides every table, applied by \texttt{src/routes/trading/control.py} only. Applied before the generator is imported, the sequence probability takes effect; imported first, it does not; the duration table never does & \textbf{not met} \\
\S7.3, \S7.4 & classifier AUC-ROC $\approx$0.993 across all 52 Thames gauges; GAUGE-8f1d368d & the training summary: 11 gauges, mean 0.9935; GAUGE-8f1d368d absent & \textbf{not met} \\
\bottomrule
\end{longtable}}


% =========================================================================
\mkmlimitations
% =========================================================================

\begin{enumerate}[leftmargin=*]
\item \textbf{Benchmark, not validation.} The code is shown to compute its documented
  equations; whether sequences like these occur over the Thames is not tested.
\item \textbf{Searches are of this repository,} at the commit above. The classifier
  evidence is the data store's training summary, retrained on v3.1's catalogue.
\end{enumerate}

\mkmhistorypage{07-Nov-2026 & 0.1 & First issue: tests the filed v2.0; 11 of its 17 claims fail & David K Kelly \\
08-Nov-2026 & 0.2 & Tests v3.0; all 17 of its claims hold & David K Kelly \\
09-Nov-2026 & 0.3 & Tests v3.1, with Tables~1 and~3 recalibrated; \bmClaimsMet{} of its \bmClaims{} claims hold & David K Kelly \\}

\end{document}
